% 教学数据由四个单位向量的内积构造，温度为0.2。
\begin{tikzpicture}[x=1cm,y=1cm,font=\figurefont\small,
  arr/.style={-{Stealth[length=2mm]},BookInk,line width=0.8pt}]
  \node[anchor=west] at (0,4.5) {缩放后分数$\ell_{i,j}=s_{i,j}/\tau$};
  \node[align=center] at (3.8,3.55) {$q_1$\\红色汽车};
  \node[align=center] at (6.0,3.55) {$q_2$\\红色自行车};
  \node[anchor=east] at (2.6,2.5) {图像1};
  \node[anchor=east] at (2.6,1.2) {图像2};
  \foreach \x/\y/\val in {3.8/2.5/4,6.0/2.5/3,3.8/1.2/3,6.0/1.2/4}{
    \draw[BookRule,line width=0.5pt,fill=BookPaper]
      (\x-1,\y-0.55) rectangle (\x+1,\y+0.55);
    \node at (\x,\y) {$\val$};
  }
  \draw[BookTeal,line width=1.2pt] (2.8,1.95) rectangle (4.8,3.05);
  \draw[BookTeal,line width=1.2pt] (5.0,0.65) rectangle (7.0,1.75);
  \draw[arr] (7.2,2.5) -- (8.15,2.5);
  \node[anchor=west,align=left] at (8.2,2.5) {行softmax\\$0.7311$\\$0.2689$};
  \draw[arr] (3.8,0.4) -- (3.8,-0.2);
  \node[align=center,anchor=north] at (3.8,-0.3)
    {列softmax\\$(0.7311,\;0.2689)^\top$};
  \node[anchor=west] at (0,-1.35) {对角线：数据配对\qquad 非对角线：批次负例};
\end{tikzpicture}
